\chapter{Planificación}

En la transformación entre la fase del análisis (donde se definen los requisitos y el
funcionamiento del sistema) y la fase de desarrollo es necesario establecer una estrategia
de ejecución detallada. El objetivo de este capítulo es establecer la metodología, la técnica
y la parte económica de esta aplicación web.

Para ello se detallan todos los recursos necesarios, organizando las tareas en el tiempo
mediante una planificación en forma de calendario, y se analizan las tecnologías elegidas
para programar y desarrollar cada punto. Para finalizar, se calcula el coste final estimado
del proyecto.

\section{Organización y recursos}

Una buena estructura a la hora de desarrollar un proyecto software como este desde el
inicio es fundamental para poder realizarlo de forma exitosa y evitando complicaciones. En
este caso, se utiliza una estructura ágil y eficaz basada en sprints (ciclos cortos de trabajo),
con el objetivo de ir creando pequeñas versiones funcionales del proyecto de manera
continua. Esto aumenta tanto la claridad del desarrollo como la capacidad de adaptación
ante cambios o errores detectados a lo largo del proceso.

De este modo, el director del trabajo puede ver los avances e ir validando paso a paso. Por
mi parte, como alumno del TFG, asumo todos los roles del equipo de desarrollo, llevando a
cabo todas las fases de este proyecto.

Además de definir la metodología, es fundamental establecer con qué medios se cuenta
para poder realizar el proyecto. Esto se divide en dos grupos principales: los recursos
humanos y los recursos materiales. A continuación se analizan en detalle estos puntos.

\subsection{Recursos humanos}

Como se ha mencionado previamente, y es lógico en un TFG, se asumen todos los roles que
se describen en este apartado, pero para poder hacer una estimación realista de los costes
se toma como referencia el mercado laboral actual de España. Se simula que este proyecto
lo llevan a cabo varios profesionales.

Dependiendo del sprint en el que se esté, hace falta un perfil técnico u otro. Teniendo esto
en cuenta, se han definido los siguientes roles, que serán necesarios a lo largo del proyecto.

\begin{table}[H]
	\centering
	\caption{Salario medio del personal.}
	\begin{tabularx}{0.75\textwidth}{X r r}
		\toprule
		\textbf{Rol} & \textbf{Salario anual bruto} & \textbf{Salario por hora} \\
		\midrule
		Analista                  & 31.500\,\euro{} & 16,15\,\euro{} \\
		Ingeniero de Datos        & 30.000\,\euro{} & 15,38\,\euro{} \\
		Ingeniero de Requisitos   & 35.000\,\euro{} & 17,95\,\euro{} \\
		Ingeniero de Software     & 35.000\,\euro{} & 17,95\,\euro{} \\
		Programador               & 28.500\,\euro{} & 14,62\,\euro{} \\
		\bottomrule
	\end{tabularx}
\end{table}

Estos datos son una aproximación basada en ofertas reales, pero pueden variar en función
del momento en el que se consulten o dependiendo del nivel y experiencia de los
profesionales.

A lo largo de cada tarea se reparten estos roles. En este caso, al ser una aplicación web de
visualización 3D, las tareas se dividen entre backend y frontend. Aunque las tareas
concretas pueden cambiar en cada ciclo, se han definido las siguientes asignaciones.

\begin{table}[H]
	\centering
	\caption{Secciones generales de cada sprint.}
	\begin{tabularx}{\textwidth}{X >{\centering\arraybackslash}p{1.8cm} l}
		\toprule
		\textbf{Sección} & \textbf{Capa} & \textbf{Encargado} \\
		\midrule
		Procesamiento de datos y algoritmos        & Backend  & Ingeniero de Datos \\
		Diseñar la API y su comunicación           & Backend  & Ingeniero de Software \\
		Diseñar el entorno y la visualización 3D   & Frontend & Ingeniero de Software \\
		Programación e integración del código      & Ambas    & Programador \\
		Pruebas de software y validación           & Ambas    & Ingeniero de Software \\
		\bottomrule
	\end{tabularx}
\end{table}

Partiendo de esta base, y sabiendo que la aplicación sigue un modelo Cliente-Servidor con
contenedores, más adelante se calcula cuánto cuesta el trabajo humano y cómo se reparten
sus horas a lo largo del tiempo de trabajo.

\subsection{Recursos materiales}

Para poder sacar adelante la aplicación también es fundamental contar con unas buenas
herramientas, además del personal técnico. Se necesita software y hardware a la altura de
este proyecto para poder programarlo y hacer pruebas de calidad previas a su salida.

Para poder obtener un presupuesto y saber cuánto cuesta cada una de estas herramientas,
que son más abstractas, se toma como referencia lo que marca la Ley del Impuesto sobre
Sociedades. Según esta ley, se puede estimar que los productos software tienen una vida
útil de 6 años (2190 días) antes de quedar obsoletos. En cuanto a las herramientas hardware,
la estimación está fijada en 8 años (2920 días).

Una ventaja que se encuentra es que la mayoría de las tecnologías utilizadas (como Docker,
VS Code, Python y React) son gratuitas o de código abierto, lo que hace el proyecto más
económico. En la siguiente tabla se puede ver todo de forma más detallada.

\begin{table}[H]
	\centering
	\caption{Productos software utilizados.}
	\begin{tabularx}{\textwidth}{X r c r}
		\toprule
		\textbf{Licencias} & \textbf{Coste total} & \textbf{Amortización} & \textbf{Coste diario} \\
		\midrule
		Microsoft Windows 11 Pro               & 12,90\,\euro{} & 2190 días & 0,005\,\euro{}/día \\
		Microsoft Office 2024 Pro              & 12,49\,\euro{} & 2190 días & 0,005\,\euro{}/día \\
		VS Code y otras tecnologías gratuitas  & 0,00\,\euro{}  & 2190 días & 0,000\,\euro{}/día \\
		\midrule
		\multicolumn{3}{r}{\textbf{Total}} & \textbf{0,010\,\euro{}/día} \\
		\bottomrule
	\end{tabularx}
\end{table}

\begin{table}[H]
	\centering
	\caption{Productos hardware utilizados.}
	\begin{tabularx}{\textwidth}{X l r c r}
		\toprule
		\textbf{Producto} & \textbf{Tipo} & \textbf{Coste total} & \textbf{Amortización} & \textbf{Coste diario} \\
		\midrule
		ASUS Zenbook 14         & Portátil & 1.099,00\,\euro{} & 2920 días & 0,38\,\euro{}/día \\
		HP OMEN 24 180Hz        & Monitor  & 169,00\,\euro{}   & 2920 días & 0,06\,\euro{}/día \\
		Logitech G203           & Ratón    & 29,99\,\euro{}    & 2920 días & 0,01\,\euro{}/día \\
		Teclado inalámbrico TKL & Teclado  & 69,99\,\euro{}    & 2920 días & 0,02\,\euro{}/día \\
		\midrule
		\multicolumn{4}{r}{\textbf{Total}} & \textbf{0,47\,\euro{}/día} \\
		\bottomrule
	\end{tabularx}
\end{table}

\section{Planificación temporal}

Para poder organizar la planificación en forma de calendario de una forma acertada, se
toman jornadas laborales estándar de 8 horas; los sprints mencionados en puntos anteriores
son la mejor forma de estructurar el trabajo en la línea temporal.

Es fundamental organizar el orden de preferencia de las tareas. Por ejemplo, no se pueden
programar los filtros visuales si aún no se ha realizado la renderización de las conexiones en
el modelo 3D. Por ello, los requisitos funcionales se clasifican en el siguiente orden de
prioridad, en una escala de 1 a 4, donde los requisitos con prioridad 1 son esenciales y
prioritarios, mientras que los de nivel 4 son utilidades secundarias. Esta sección se refiere
únicamente al desarrollo de software puro.

\begin{table}[H]
	\centering
	\caption{Priorización de requisitos.}
	\begin{tabularx}{\textwidth}{>{\centering\arraybackslash}p{1.3cm} X >{\centering\arraybackslash}p{1.8cm}}
		\toprule
		\textbf{RF} & \textbf{Resumen de funcionalidad} & \textbf{Prioridad} \\
		\midrule
		RF 1  & Carga de red cerebral (ficheros locales o atlas precargados). & 1 \\
		RF 2  & Renderizado espacial 3D de nodos y aristas ponderadas. & 1 \\
		RF 3  & Interacción y controles de cámara orbital (rotar, zoom, paneo). & 1 \\
		RF 4  & Filtrado dinámico de conexiones visibles mediante umbral. & 2 \\
		RF 5  & Consulta de información detallada y métricas de cada nodo. & 2 \\
		RF 6  & Selección e inspección visual de un nodo haciendo clic. & 2 \\
		RF 7  & Personalización gráfica (colores/tamaños según atributos biológicos). & 3 \\
		RF 8  & Panel lateral (HUD) con estadísticas globales de la red activa. & 3 \\
		RF 9  & Exportación y descarga de la visualización actual como imagen. & 4 \\
		RF 10 & Restablecimiento rápido de la vista y parámetros iniciales. & 4 \\
		\bottomrule
	\end{tabularx}
\end{table}

Una vez claros los requisitos más urgentes, se pueden asignar a cada uno de los sprints. En
total, se divide el desarrollo de la plataforma en cuatro sprints. En la siguiente tabla se
muestra cómo quedan distribuidas estas funcionalidades. Junto a cada bloque se incluye la
estimación de días laborables que llevará completarlo, medición fundamental para montar
el calendario definitivo.

\begin{table}[H]
	\centering
	\caption{Asignación de requisitos por sprint.}
	\begin{tabularx}{\textwidth}{>{\centering\arraybackslash}p{1.6cm} >{\centering\arraybackslash}p{2.6cm} >{\raggedright\arraybackslash}X}
		\toprule
		\textbf{Sprint} & \textbf{Estimación temporal} & \textbf{Requisitos funcionales asignados} \\
		\midrule
		1 & 18 días &
			RF 1. Carga de red cerebral (ficheros locales o atlas precargados). \newline
			RF 2. Renderizado espacial 3D de nodos y aristas ponderadas. \\
		\addlinespace
		2 & 13 días &
			RF 3. Interacción y controles de cámara orbital (rotar, zoom...). \newline
			RF 4. Filtrado dinámico de conexiones visibles mediante umbral. \\
		\addlinespace
		3 & 12 días &
			RF 5. Consulta de información detallada y métricas de cada nodo. \newline
			RF 6. Selección e inspección visual de un nodo haciendo clic. \\
		\addlinespace
		4 & 8 días &
			RF 7. Personalización gráfica (colores/tamaños según atributos biológicos). \newline
			RF 8. Panel lateral (HUD) con estadísticas globales de la red activa. \newline
			RF 9. Exportación y descarga de la visualización actual como imagen. \newline
			RF 10. Restablecimiento rápido de la vista y parámetros iniciales. \\
		\bottomrule
	\end{tabularx}
\end{table}

\section{Análisis de tecnologías disponibles}

El objetivo de este proyecto es el diseño, desarrollo y despliegue de una aplicación web
interactiva de visualización 3D de redes neuronales. La estructura de esta aplicación se
divide en dos componentes principales:

\begin{itemize}
	\item \textbf{Servicio de backend:} encargado de llevar a cabo el procesamiento y análisis
		de las redes neuronales. Su función es aplicar análisis a los grafos de las fuentes de
		datos y exponer los resultados a través de una API.
	\item \textbf{Aplicación cliente o frontend:} encargada de la interacción con el usuario y de
		la visualización 3D en el navegador del cliente.
\end{itemize}

La selección de estas tecnologías se ha basado en la idoneidad para el problema específico.
A continuación se analizan y comparan las tecnologías consideradas.

\subsection{Framework de Backend y Análisis de Grafos}

Para el backend se debe tener en cuenta que debe cumplir con dos requisitos clave:

\begin{itemize}
	\item Servir los datos de las redes al cliente a través de una API.
	\item Utilizar bibliotecas de análisis de grafos para el procesamiento de las matrices de
		adyacencia.
\end{itemize}

Ya que Python contiene bibliotecas de análisis de grafos implementadas de forma nativa
(como NetworkX, igraph o graph-tool), elegir un framework de backend en Python es la
decisión más eficiente y lógica. Con esto se evita la dificultad de crear servicios de
intercomunicación entre lenguajes.

\begin{table}[H]
	\centering
	\caption{Alternativas de backend en Python.}
	\small
	\begin{tabularx}{\textwidth}{>{\raggedright\arraybackslash}p{1.8cm} >{\raggedright\arraybackslash}p{3.2cm} >{\raggedright\arraybackslash}X}
		\toprule
		\textbf{Framework} & \textbf{Arquitectura} & \textbf{Utilidad para el proyecto} \\
		\midrule
		Django & Monolítico (Full-stack, ``pilas incluidas'') &
			Incluye un ORM, un panel de administrador y un sistema de autenticación.
			Estos elementos no son imprescindibles para el núcleo de esta aplicación web,
			la cual se enfoca en un servicio de datos. \\
		\addlinespace
		Flask & Micro-framework (Minimalista) &
			Para desarrollar APIs REST livianas, Flask es perfecto. Posibilita una
			integración simple y directa con librerías como NetworkX. Su filosofía
			minimalista posibilita construir únicamente lo que se necesita, lo cual es
			ideal para un servicio de API limitado. \\
		\addlinespace
		FastAPI & Micro-framework (Moderno, Asíncrono) &
			Proporciona un rendimiento más elevado debido a su carácter asíncrono
			(ASGI). Produce de manera automática documentación interactiva
			(OpenAPI/Swagger), lo que representa un beneficio importante para el
			desarrollo y la depuración de la API. \\
		\bottomrule
	\end{tabularx}
\end{table}

\textbf{Selección para el backend:}

Django es muy complejo y su enfoque monolítico no ayuda al desarrollo de la aplicación,
por lo que se descarta, quedando las opciones de Flask y FastAPI:

\begin{itemize}
	\item \textbf{Flask:} permite tener un endpoint funcional de forma rápida y fácil de
		integrar. Es la opción más directa, siendo sencilla también la integración con
		NetworkX.
	\item \textbf{FastAPI:} tiene un mayor rendimiento asíncrono, es una alternativa más
		moderna y permite la validación de datos mediante \textit{type hints}. La
		documentación automática también es una ventaja notable.
\end{itemize}

\textbf{Selección final:} la mejor opción para desarrollar este proyecto es Flask, ya que
permite crear una API totalmente funcional para leer los ficheros de forma directa, para
después procesarlos con NetworkX y enviarlos como JSON al frontend.

\subsection{Framework de Frontend (UI)}

El frontend es la parte con la que interactúa el usuario. La aplicación web debe gestionar la
visualización 3D y controlar el estado de la aplicación.

\begin{table}[H]
	\centering
	\caption{Alternativas de frontend.}
	\small
	\begin{tabularx}{\textwidth}{>{\raggedright\arraybackslash}p{2.1cm} >{\raggedright\arraybackslash}p{3.0cm} >{\raggedright\arraybackslash}p{3.0cm} >{\raggedright\arraybackslash}X}
		\toprule
		\textbf{Framework} & \textbf{Paradigma} & \textbf{Ecosistema} & \textbf{Integración 3D} \\
		\midrule
		React & Basado en componentes (Declarativo) &
			Enorme. Mantenido por Meta. &
			Contiene bindings que integran Three.js. \\
		\addlinespace
		Vue.js & Basado en componentes (Progresivo) &
			Grande. Impulsado por la comunidad. &
			Existen integraciones como \textit{trois}, que también integran Three.js. \\
		\addlinespace
		Angular & Framework completo (MVC / Componentes) &
			Grande. Mantenido por Google. &
			La integración con librerías 3D es menos común o de más bajo nivel. \\
		\bottomrule
	\end{tabularx}
\end{table}

\textbf{Selección para el frontend:} para decidir entre estas tres opciones:

\begin{itemize}
	\item \textbf{Angular:} no aporta beneficios respecto al resto y queda descartado por su
		dificultad y rigidez.
	\item \textbf{React y Vue:} son mejores opciones que Angular, pero React tiene una ventaja
		superior gracias al ecosistema de visualización 3D, en concreto al binding
		\textit{react-three-fiber}, que permite desarrollar con bastante potencia el
		renderizado 3D usando Three.js.
\end{itemize}

\textbf{Selección final:} la mejor opción para el desarrollo de la aplicación web a nivel de UI
es React, que por su gran estabilidad, popularidad y sobre todo su buena integración 3D es
la mejor forma de gestionar el desafío de una interfaz 3D interactiva.

\subsection{Tecnología de visualización 3D}

La biblioteca elegida para la visualización 3D debe ser capaz de renderizar grafos en un
canvas WebGL y permitir una interacción fluida para una buena experiencia de usuario.

\begin{table}[H]
	\centering
	\caption{Tecnologías de visualización 3D consideradas.}
	\small
	\begin{tabularx}{\textwidth}{>{\raggedright\arraybackslash}p{2.0cm} >{\raggedright\arraybackslash}p{2.4cm} >{\raggedright\arraybackslash}p{2.1cm} >{\raggedright\arraybackslash}X}
		\toprule
		\textbf{Tecnología} & \textbf{Tipo} & \textbf{Nivel de abstracción} & \textbf{Caso de uso} \\
		\midrule
		PlayCanvas & Motor de juego 3D (con editor) & Alto &
			Motor enfocado al desarrollo de videojuegos, con editor visual en la nube,
			pero su modelo no se adapta a la carga dinámica de datos. \\
		\addlinespace
		luma.gl / deck.gl & Framework de visualización de datos & Alto (geodatos) &
			Optimizado para la carga y visualización de grandes cantidades de datos
			geoespaciales, como mapas. No diseñado para grafos abstractos en 3D. \\
		\addlinespace
		graph.gl & Framework de visualización de grafos & Alto (grafos) &
			Proyecto con una comunidad pequeña y menos mantenido que el resto de
			alternativas. Alto riesgo de encontrar problemas y poca documentación. \\
		\addlinespace
		Babylon.js & Motor de juego / framework 3D & Medio-Alto &
			Mantenido por Microsoft. Es un ``todo en uno'' muy potente y
			relativamente sencillo. Ecosistema robusto. \\
		\addlinespace
		Three.js & Librería de renderizado 3D & Bajo &
			Es el estándar en la industria. La librería WebGL de más bajo nivel y más
			utilizada. Requiere más esfuerzo en la gestión de escena, cámaras, luces y
			bucle de renderizado manual. \\
		\bottomrule
	\end{tabularx}
\end{table}

\textbf{Selección para la visualización 3D:} se descartan PlayCanvas, luma.gl y graph.gl por
no ajustarse al uso de visualización de datos 3D. La mejor elección se encuentra entre:

\begin{itemize}
	\item \textbf{Babylon.js:} una excelente opción ``todo en uno'' y relativamente sencilla de
		implementar.
	\item \textbf{Three.js:} más flexible y con más ejemplos disponibles.
\end{itemize}

Esta decisión está ligada a la elección previa del framework de frontend. Al haberse
seleccionado React por sus ventajas respecto al resto de opciones, lo más lógico es usar
Three.js, aunque no de forma pura, sino a través de un render declarativo.

\textbf{Selección final:} se utiliza Three.js, mediante \textit{react-three-fiber}. Esta
combinación mejora el equilibrio del proyecto gracias a:

\begin{itemize}
	\item \textbf{Potencia:} Three.js es la librería más potente y flexible.
	\item \textbf{Facilidad de uso:} al usar React para gestionar la aplicación, la
		implementación de Three.js es más sencilla gracias a react-three-fiber.
	\item \textbf{Integración en el sistema:} react-three-fiber actúa como la conexión perfecta
		entre ambos, permitiendo declarar componentes como luces o cámaras igual que
		componentes de React, eliminando así las dificultades del bucle de renderizado
		manual y facilitando la conexión con el estado de la aplicación.
\end{itemize}

\subsection{Resumen del stack seleccionado}

\begin{table}[H]
	\centering
	\caption{Resumen del stack tecnológico seleccionado.}
	\small
	\begin{tabularx}{\textwidth}{>{\raggedright\arraybackslash}p{2.2cm} >{\raggedright\arraybackslash}p{2.3cm} >{\raggedright\arraybackslash}p{2.9cm} >{\raggedright\arraybackslash}X}
		\toprule
		\textbf{Componente} & \textbf{Tecnología principal} & \textbf{Bibliotecas} & \textbf{Justificación} \\
		\midrule
		Backend & Python & Flask, NetworkX &
			Integración nativa para análisis de grafos, eficiencia y enfoque minimalista
			para APIs REST. \\
		\addlinespace
		Frontend & JavaScript & React, Three.js, react-three-fiber &
			Modelo declarativo de componentes, ecosistema líder y la mejor
			integración de su clase entre UI y renderizado 3D. \\
		\addlinespace
		Despliegue & Docker & Docker, Docker Compose &
			Cumple con los requisitos del proyecto, garantiza la reproducibilidad del
			entorno y facilita la orquestación de los servicios de backend y frontend. \\
		\bottomrule
	\end{tabularx}
\end{table}

\section{Estimación de costes}

Para cerrar el capítulo de la planificación, se recopilan todos los datos calculados en los
apartados anteriores y se presenta la estimación de costes final del proyecto.

Como se ha visto en el resumen de las tareas, el esfuerzo en recursos humanos marca que el
proyecto tiene una duración total de 62 días laborables. Este es exactamente el periodo de
tiempo utilizado para calcular la parte proporcional de desgaste (amortización) que sufren
los equipos físicos y las licencias de sistema operativo durante el desarrollo.

\begin{table}[H]
	\centering
	\caption{Estimación de costes total del proyecto.}
	\begin{tabularx}{0.85\textwidth}{X r}
		\toprule
		\textbf{Recursos} & \textbf{Coste} \\
		\midrule
		Recursos humanos                                        & 8.010,80\,\euro{} \\
		Recursos hardware (62 días $\times$ 0,47\,\euro{}/día)  & 29,14\,\euro{} \\
		Recursos software (62 días $\times$ 0,014\,\euro{}/día) & 0,87\,\euro{} \\
		Tecnologías utilizadas                                  & 0,00\,\euro{} \\
		\midrule
		\textbf{Coste total} & \textbf{8.040,81\,\euro{}} \\
		\bottomrule
	\end{tabularx}
\end{table}
